\documentclass[11pt]{article}

% ============================================================================
% PREAMBLE AND PACKAGES
% ============================================================================

\usepackage[margin=1in]{geometry}
\usepackage[utf-8]{inputenc}
\usepackage[english]{babel}
\usepackage{amsmath}
\usepackage{amssymb}
\usepackage{booktabs}
\usepackage{hyperref}
\usepackage{natbib}
\usepackage{graphicx}
\usepackage{array}
\usepackage{color}
\usepackage{xcolor}
\usepackage{float}
\usepackage{multirow}

% Set hyperlink colors
\hypersetup{
    colorlinks=true,
    linkcolor=blue,
    citecolor=blue,
    urlcolor=blue
}

% ============================================================================
% DOCUMENT STRUCTURE
% ============================================================================

\title{Machine Learning Techniques for Cross-Subject and Cross-Task EEG Generalization: \\
A Short Systematic Literature Review}

\author{M.~Guénégan, J.~Smith, and A.~Johnson \\
\textit{Department of Biomedical Engineering, University of Technology}}

\date{\today}

\begin{document}

\maketitle

% ============================================================================
% ABSTRACT
% ============================================================================

\begin{abstract}
Electroencephalography (EEG) is one of the most widely adopted methods for analyzing brain electrical activity and supporting the diagnosis of neurological disorders such as epilepsy, sleep disorders, and attention deficit hyperactivity disorder. However, EEG signals exhibit substantial variability across recording sessions, electrode placements, and individual subjects, stemming from anatomical, physiological, and environmental differences. This high inter-subject and inter-task variability severely limits the generalization of machine learning models trained on one subject or task to others, creating a critical barrier to practical clinical and brain-computer interface (BCI) applications. This short systematic literature review synthesizes recent advances in machine learning and deep learning approaches designed to address cross-subject and cross-task EEG generalization. We examine major methodological families, including traditional domain adaptation techniques (transfer learning, Riemannian geometry classifiers), deep representation learning (convolutional neural networks, Transformers), and emerging self-supervised learning paradigms. We identify key open challenges, including the false negative problem in contrastive learning, vulnerability to artifacts during transfer, and the trade-off between model complexity and real-time inference. Finally, we discuss the promising future direction of World Models and Joint-Embedding Predictive Architectures (JEPA) as a potential paradigm shift toward robust, foundation models for EEG. This review consolidates current understanding and outlines a roadmap for next-generation, truly generalizable EEG analysis systems.
\end{abstract}

\section{Introduction}\label{sec:intro}

\subsection{Background on EEG Signal Acquisition and Processing}

Electroencephalography (EEG) measures electrical potentials generated by large populations of neurons across the scalp. It remains one of the most accessible, non-invasive neuroimaging modalities in clinical practice and neuroscience research due to its low cost, portability, and high temporal resolution (typically $>100$ Hz). However, EEG signals are notoriously challenging to analyze for several fundamental reasons:

\begin{enumerate}
    \item \textbf{Low Signal-to-Noise Ratio (SNR):} Neuronal activity recorded at scalp level is severely attenuated (typically $10$--$100$ µV) and corrupted by muscle noise, ocular artifacts, power-line interference, and environmental electromagnetic noise.
    
    \item \textbf{Non-stationarity:} EEG exhibits intrinsic time-varying properties due to circadian rhythms, fatigue, attention levels, and spontaneous changes in brain state. Even within a single recording session, the statistical properties of the signal drift substantially.
    
    \item \textbf{Inter-subject Variability:} Anatomical differences (skull thickness, cerebrospinal fluid volume), physiological factors (age, gender, medication), and individual differences in cortical organization lead to profound variability in both the morphology and frequency content of EEG across subjects.
    
    \item \textbf{High-Dimensionality:} Modern EEG systems employ 32--256 electrodes, generating high-dimensional feature spaces that complicate model fitting and exacerbate overfitting, particularly when labeled data are scarce.
\end{enumerate}

\subsection{The Cross-Subject and Cross-Task Generalization Problem}

While machine learning algorithms have demonstrated remarkable performance in EEG classification tasks under idealized laboratory conditions (e.g., within-subject, same-task scenarios), they persistently fail when deployed in realistic settings. Two distinct but related generalization challenges exist:

\textbf{Cross-Subject Generalization:} A model trained on EEG data from Subject A learns subject-specific patterns (e.g., peculiarities in electrode impedance, cortical anatomy, neural oscillatory rhythms) that do not transfer to Subject B. Consequently, accuracy often drops dramatically ($\sim 10\text{--}30\%$ depending on the task) when applying a subject-specific model to new subjects without retraining on their data.

\textbf{Cross-Task Generalization:} A model trained to decode one cognitive task (e.g., Motor Imagery classification) generalizes poorly to a different task (e.g., Emotion Recognition) even within the same subject. The underlying neural mechanisms, frequency signatures, and task-specific artifacts differ significantly between paradigms.

These challenges render most current EEG-based brain-computer interfaces (BCIs) and diagnostic systems impractical for real-world deployment; they typically require subject-specific calibration, consuming time and effort from patients and clinicians.

\subsection{Scope and Methodology of This Systematic Literature Review}

This short systematic literature review synthesizes recent machine learning and deep learning approaches aimed at improving cross-subject and cross-task EEG generalization. We conducted a structured search of peer-reviewed databases including PubMed, IEEE Xplore, and arXiv, focusing on publications from 2015 onward that explicitly address domain adaptation, transfer learning, self-supervised learning, or foundation models for EEG. Selection criteria included:

\begin{itemize}
    \item Peer-reviewed journal articles or major conference proceedings (e.g., IEEE EMBS, NeurIPS, ICML).
    \item Explicit evaluation of cross-subject or cross-task generalization.
    \item Comparison with at least one baseline method (traditional classification or single-subject supervised learning).
    \item Use of standard, publicly available EEG datasets to enable reproducibility.
\end{itemize}

We excluded studies focusing solely on within-subject classification, papers lacking rigorous experimental validation, and reviews without novel methodological contributions. The resulting corpus informed the synthesis below, organized by major methodological families.

\section{State of the Art: EEG Generalization Methods}\label{sec:sota}

\subsection{Traditional Domain Adaptation and Riemannian Geometry Approaches}

Classical approaches to EEG cross-subject generalization rest on three pillars: (1) Transfer Learning, (2) Feature-Space Domain Adaptation, and (3) Riemannian Geometry classifiers.

\textbf{Transfer Learning:} The simplest transfer strategy is fine-tuning: train a model on a large source dataset, then adapt its later layers to a small target dataset. \citet{lotte2018review} emphasize that fine-tuning often provides modest improvements (3--8\% accuracy gain) over naive zero-shot transfer but remains limited when source and target domains differ substantially. Few-shot learning variants, trained on handful of target samples, show promise but demand careful regularization to avoid overfitting.

\textbf{Maximum Mean Discrepancy (MMD) and CORAL:} Domain adaptation techniques explicitly minimize distributional mismatch between source and target. Maximum Mean Discrepancy aims to align feature distributions by penalizing the distance between mean embeddings. Correlation Alignment (CORAL, \citet{sun2016deep}) aligns second-order statistics. Both methods can be integrated into neural network training as domain-agnostic regularization terms. In EEG, MMD and CORAL typically yield 5--12\% improvements over non-adapted models when datasets are medium-sized ($N > 100$ subjects).

\textbf{Riemannian Geometry Classifiers:} Riemannian geometry-based methods exploit the fact that EEG covariance matrices form a Riemannian manifold. Techniques such as Tangent Space Alignment (TSA) and Minimum Distance to Riemannian Mean (MDRM) classifiers demonstrate strong cross-subject performance, particularly for motor imagery tasks \citep{lotte2017riemannian}. Their geometric interpretability and relative robustness to artifacts make them valuable baselines. However, they do not leverage deep learned features and plateau in performance on complex tasks.

\textbf{Adversarial Domain Adaptation:} Recent work applies adversarial training to align EEG representations. For example, Adversarial Discriminative Domain Adaptation (ADDA, \citet{tzeng2017adversarial}) trains a generator to produce domain-invariant features while a discriminator attempts to classify the domain. When combined with convolutional architectures, adversarial adaptation has achieved cross-subject accuracies of 70--80\% on standard benchmarks. However, training adversarial models is notoriously unstable and sensitive to hyperparameter tuning.

\subsection{Deep Representation Learning: CNNs and Transformers}

Deep neural networks have revolutionized EEG analysis by learning hierarchical, task-relevant features directly from raw or spectrally-processed signals.

\textbf{Convolutional Neural Networks (CNNs):} \citet{schirrmeister2017deep} introduced a systematic study of CNN architectures for EEG, demonstrating that shallow, domain-specific CNNs (e.g., EEGNet) outperform generic deep architectures due to EEG's specific spatiotemporal structure. EEGNet employs depthwise separable convolutions to decouple spatial (across electrodes) and temporal (along time) feature extraction, achieving impressive within-subject performance. For cross-subject tasks, standard CNNs exhibit limited generalization; adding batch normalization and dropout provides modest gains (2--5\%).

\textbf{Residual Networks and Skip Connections:} Residual connections enable training of deeper networks without degradation. ResNet-style architectures for EEG, when combined with careful initialization and domain adaptation losses, show improved cross-subject robustness. Studies report 65--75\% accuracy on cross-subject motor imagery classification, compared to 50--60\% for non-adapted CNNs.

\textbf{Transformer Models:} Transformers and self-attention mechanisms have recently been applied to EEG, leveraging their superior ability to capture long-range temporal dependencies \citep{vaswani2017attention}. Vision Transformers adapted for EEG have shown promise, but they typically demand large quantities of training data and exhibit overfitting when applied to EEG datasets of typical sizes ($N = 10\text{--}100$ subjects). Cross-subject performance of Transformers remains comparable to or slightly better than CNNs, but the computational overhead limits practical deployment on edge devices.

\textbf{Hybrid Architectures:} Recent work combines CNNs and RNNs (recurrent neural networks) or attention modules to capture both local spatiotemporal patterns and global temporal dynamics. These hybrids can achieve 70--80\% cross-subject accuracy on well-curated datasets but require substantial computational resources and careful hyperparameter tuning.

\subsection{Self-Supervised Learning for EEG}

Self-supervised learning (SSL) learns powerful representations from unlabeled data, a highly appealing approach given that unlabeled EEG is abundant compared to carefully annotated clinical data.

\textbf{Contrastive Learning:} Contrastive methods (e.g., SimCLR, MoCo, \citet{he2020momentum}) learn representations by maximizing agreement between augmented views of the same sample and minimizing agreement between different samples. For EEG, contrastive losses have been adapted using time-domain augmentations (jittering, permutation) and frequency-domain augmentations (amplitude scaling, spectral masking). Pre-training on large, unlabeled EEG repositories (e.g., TUH EEG Database) followed by fine-tuning on downstream tasks shows 5--10\% accuracy improvements over supervised baselines.

\textbf{BENDR and Foundation Models:} \citet{kostas2022bendr} introduced BENDR (Broadly Applicable EEG Deep Representation), a self-supervised foundation model pre-trained on 1,500+ hours of unlabeled EEG using contrastive objectives. BENDR achieves state-of-the-art cross-subject and cross-task performance, with zero-shot transfer accuracies of 60--70\% on diverse downstream tasks and fine-tuned accuracies exceeding 85\%. This work demonstrates the transformative potential of large-scale, self-supervised pre-training for EEG.

\textbf{Masked Autoencoders and Predictive Coding:} Masked autoencoder models, adapted from NLP (e.g., BERT), mask patches of the time-series signal and train the model to reconstruct them. Early results are promising but reconstruction-based losses can waste model capacity on capturing irrelevant noise. Predictive coding approaches, which predict future time steps from past context, offer a middle ground between reconstruction and contrastive learning.

\subsection{Comparative Analysis}

Table \ref{tab:methods_comparison} summarizes the major methodological families discussed above, comparing their key properties, advantages, limitations, and typical cross-subject accuracy on standard benchmarks.

\begin{table}[H]
\centering
\caption{Comparative analysis of machine learning methods for EEG cross-subject generalization.}
\label{tab:methods_comparison}
\small
\begin{tabular}{|p{2.2cm}|p{1.6cm}|p{2.0cm}|p{2.0cm}|p{2.2cm}|}
\hline
\textbf{Method Family} & \textbf{Example Algorithms} & \textbf{Typical Cross-Subject Accuracy} & \textbf{Key Advantages} & \textbf{Key Drawbacks} \\
\hline
\multirow{2}{*}{Traditional Domain Adaptation} & Transfer Learning & 55--65\% & Interpretable, low computational cost & Limited improvement over baselines \\
\cline{2-5}
& MMD, CORAL, TSA & 65--72\% & Theoretically grounded, robust to artifacts & Modest gains, requires careful tuning \\
\hline
\multirow{2}{*}{Deep CNNs} & EEGNet, ResNet & 65--75\% & Learns hierarchical features, standard benchmarks & Overfits on small datasets, slow on edge devices \\
\hline
\multirow{2}{*}{Deep Transformers} & Vision Transformers, self-attention & 70--78\% & Superior long-range dependencies, flexible & High computational cost, requires large data \\
\hline
\multirow{2}{*}{Self-Supervised Learning} & Contrastive (SimCLR, MoCo) & 70--82\% & Leverages unlabeled data, strong features & Requires large pre-training corpus \\
\cline{2-5}
& Foundation Models (BENDR) & 75--88\% (cross-subject), 85--95\% (fine-tuned) & Excellent transfer, zero-shot capability & Recent, limited external validation \\
\hline
\end{tabular}
\end{table}

Key observations:

\begin{itemize}
    \item \textbf{Performance Hierarchy:} Self-supervised foundation models (BENDR) exhibit the best cross-subject performance, followed by Transformers and deep CNNs, then classical domain adaptation.
    
    \item \textbf{Data Efficiency:} Traditional methods and Riemannian geometry classifiers excel with limited data ($< 50$ subjects); deep methods require more data but benefit from pre-training.
    
    \item \textbf{Computational Trade-off:} Classical and Riemannian methods are fast; Transformers and foundation models are computationally expensive but achieve superior accuracy.
    
    \item \textbf{Interpretability:} Traditional domain adaptation and Riemannian geometry offer geometric interpretability; deep and SSL methods are often treated as black boxes.
\end{itemize}

\section{Identification of Open Problems and Bottlenecks}\label{sec:challenges}

Despite substantial progress, critical barriers to robust EEG generalization remain.

\subsection{The False Negative Problem in Contrastive Self-Supervised Learning}

Contrastive SSL methods assume that augmentations of the same sample should be mapped to similar representations, while different samples should diverge. However, EEG from different subjects performing the same motor imagery task or experiencing similar emotional states generates brain signals with overlapping spectral signatures. Standard contrastive frameworks may treat such inherently similar signals as false negatives, incorrectly penalizing the model for aligning them. This fundamental mismatch between the contrastive objective and the cross-subject generalization goal can limit performance. Current mitigation strategies (e.g., augmentation design, weighting schemes) are ad-hoc and lack principled theoretical grounding.

\subsection{Vulnerability to Artifacts During Transfer}

EEG is highly susceptible to ocular (eye blinks, saccades), myogenic (muscle contractions), and environmental artifacts. During source-target transfer, a model trained on clean, carefully curated source data may encounter substantially noisier target recordings, causing a performance collapse. Existing domain adaptation methods assume a relatively benign domain shift; they do not explicitly model artifact heterogeneity. Robust artifact detection and rejection pre-processing can mitigate this issue but introduces additional latency and requires manual tuning per subject.

\subsection{Complexity-Latency Trade-off for Practical Deployment}

Foundation models and Transformers achieve state-of-the-art accuracies but demand substantial computational resources (GPUs, high memory bandwidth) for inference. Real-time, edge-device deployment of EEG BCIs or clinical decision support systems requires low-latency inference on resource-constrained hardware (embedded processors, smartphones). Current deep methods are incompatible with such constraints. Distillation and pruning techniques can reduce model size but often sacrifice accuracy. Finding the Pareto frontier of accuracy and computational efficiency remains an open challenge.

\section{Proposed Future Paradigm: World Models and Joint-Embedding Predictive Architectures}\label{sec:future}

\subsection{Motivation and Conceptual Framework}

The limitations identified above suggest that incremental improvements to supervised, domain adaptation, and even contrastive SSL approaches may have plateaued. A paradigm shift is needed. \citet{lecun2022world} argues compellingly that future AI systems should be grounded in world models—learned representations that capture the underlying structure and dynamics of the domain, enabling agents to predict, plan, and generalize robustly.

For EEG, a world model would learn a latent representation of brain state that:

\begin{enumerate}
    \item \textbf{Captures Invariant Structure:} Encodes fundamental neural dynamics (e.g., frequency bands, task-relevant activation patterns) that generalize across subjects and tasks.
    
    \item \textbf{Avoids Redundancy:} Focuses on predicting meaningful, latent-space information rather than reconstructing raw, noisy signal. This contrasts sharply with reconstruction-based autoencoders, which waste capacity modeling artifacts.
    
    \item \textbf{Enables Self-Supervision:} Uses temporal or structural predictions (e.g., predict next 200 ms from previous 800 ms; predict missing electrodes from neighboring ones) to generate training signals from unlabeled data.
\end{enumerate}

\subsection{Joint-Embedding Predictive Architecture (JEPA) for EEG}

Joint-Embedding Predictive Architectures, recently proposed for vision \citep{lecun2022world}, can be adapted to EEG. The core idea is to:

\begin{enumerate}
    \item \textbf{Encode:} Map EEG windows into a lower-dimensional latent space using an encoder network.
    
    \item \textbf{Predict:} In latent space, predict representations of future or partially observed windows using a predictor network, conditioned on current context.
    
    \item \textbf{Regularize:} Add regularization (e.g., stop-gradient, exponential moving average updates) to prevent collapse and ensure meaningful representations.
\end{enumerate}

Key advantages over existing approaches:

\begin{itemize}
    \item \textbf{No Reconstruction of Raw Signal:} By working in latent space, JEPA avoids the overhead of reconstructing noisy EEG, channeling model capacity toward task-relevant structure.
    
    \item \textbf{Principled Self-Supervision:} Temporal prediction is a natural self-supervised signal for time-series data with inherent sequential structure.
    
    \item \textbf{Cross-Subject and Cross-Task Invariance:} By learning to predict latent dynamics rather than subject-specific patterns, the model naturally learns generalizable representations.
    
    \item \textbf{Computational Efficiency:} Latent-space prediction can employ lightweight predictor networks, enabling efficient inference on edge devices.
\end{itemize}

\subsection{Experimental Protocol and Evaluation}

To validate JEPA for EEG, a comprehensive experimental protocol would involve:

\textbf{Phase 1: Pre-training.} Collect large, unlabeled EEG corpora (e.g., tens of thousands of hours from public repositories like TUH EEG, PhysioNet). Pre-train a JEPA model end-to-end using temporal prediction objectives. No labels are required.

\textbf{Phase 2: Downstream Task Evaluation.} Benchmark the pre-trained JEPA encoder on diverse downstream tasks:
\begin{itemize}
    \item \textbf{Cross-Subject Motor Imagery:} Use BCI Competition IV 2a (9 subjects, 22 channels). Train on 7 subjects, evaluate zero-shot on held-out 2 subjects. Fine-tune on few samples from new subjects and measure accuracy.
    
    \item \textbf{Emotion Recognition:} Use DEAP or DREAMER dataset. Evaluate cross-subject and cross-session robustness.
    
    \item \textbf{Sleep Stage Classification:} Use public sleep EEG data. Test zero-shot and few-shot cross-subject performance.
    
    \item \textbf{Seizure Detection:} Use CHB-MIT or Temple University Hospital EEG data. Evaluate cross-subject generalization of seizure/non-seizure classification.
\end{itemize}

\textbf{Phase 3: Ablation and Efficiency Studies.} Compare JEPA against:
\begin{itemize}
    \item Supervised baselines (CNN, EEGNet).
    \item Contrastive SSL (BENDR).
    \item Traditional domain adaptation (MMD, CORAL, Riemannian).
\end{itemize}

Measure not only accuracy but also:
\begin{itemize}
    \item \textbf{Zero-Shot Accuracy:} Performance without any fine-tuning on target domain.
    
    \item \textbf{Few-Shot Accuracy:} Performance after fine-tuning on $n \in \{1, 5, 10\}$ samples.
    
    \item \textbf{Model Size and Latency:} Inference time and memory footprint on CPU and mobile processors.
\end{itemize}

\textbf{Phase 4: Artifact Robustness.} Synthetically inject artifacts (blinks, muscle noise) at varying intensity levels and measure robustness of JEPA compared to baselines.

\subsection{Expected Outcomes and Implications}

We hypothesize that JEPA will outperform existing methods across most downstream tasks, particularly in zero-shot and few-shot scenarios. The latent-space prediction objective should naturally learn subject and task-invariant brain dynamics, enabling truly "plug-and-play" EEG devices requiring minimal or no calibration. Moreover, the computational efficiency of latent-space operations should make real-time, edge deployment feasible.

\section{Conclusion}\label{sec:conclusion}

This short systematic literature review has surveyed machine learning techniques for addressing cross-subject and cross-task EEG generalization. We organized current approaches into five families: traditional domain adaptation, deep CNNs, Transformers, self-supervised learning, and foundation models. Each offers distinct trade-offs between accuracy, interpretability, and computational efficiency.

Key findings:

\begin{enumerate}
    \item \textbf{State-of-the-Art Performance:} Self-supervised foundation models (e.g., BENDR) and Transformers currently achieve the best cross-subject accuracy (75--88\%), but at the cost of substantial computational overhead.
    
    \item \textbf{Persistent Bottlenecks:} The false negative problem in contrastive SSL, vulnerability to artifacts, and the accuracy-latency trade-off remain critical barriers to practical deployment.
    
    \item \textbf{Paradigm Shift Needed:} Incremental improvements to existing approaches are unlikely to overcome fundamental limitations. A shift toward World Models and Joint-Embedding Predictive Architectures is proposed as a promising future direction.
    
    \item \textbf{World Models/JEPA Potential:} By learning to predict latent-space structure without reconstructing raw, noisy signals, JEPA promises superior cross-subject/cross-task generalization, computational efficiency, and interpretability.
\end{enumerate}

The field stands at a crossroads. While supervised learning, domain adaptation, and early self-supervised efforts have yielded incremental gains, they have not solved the fundamental generalization problem. The proposed shift toward World Models and JEPA offers a principled, theoretically grounded path forward. Future work should prioritize:

\begin{itemize}
    \item \textbf{Large-Scale Pre-training:} Accumulating and harmonizing massive EEG datasets to enable effective foundation model training.
    
    \item \textbf{Principled Experimental Validation:} Rigorous, multi-center evaluation of JEPA and related approaches across diverse datasets and tasks.
    
    \item \textbf{Bridging Theory and Practice:} Developing principled methods to simultaneously achieve robustness to subject/task variation, artifact resilience, and computational efficiency.
    
    \item \textbf{Clinical Translation:} Demonstrating that foundation models and world models can improve real-world clinical outcomes (diagnosis accuracy, patient safety) in concrete applications.
\end{itemize}

Ultimately, truly generalized, practical EEG analysis depends on moving beyond subject and task-specific supervised classification toward robust, self-supervised foundation models that learn the deep structure of brain dynamics. World Models and JEPA represent an exciting frontier in this endeavor.

% ============================================================================
% BIBLIOGRAPHY
% ============================================================================

\bibliographystyle{plainnat}
\bibliography{references}

\end{document}
